\subsection{单模 V(m)}

设 \((u,v)\) 是 \(k^2\) 的典则基。对 \(\mathfrak{sl}(2,k)\) 的恒同表示，

\[
X _ {+} u = 0 \quad H u = u \quad X _ {-} u = - v
\]

\[
X _ {+} v = u \quad H v = - v \quad X _ {-} v = 0.
\]

考察对称代数 \(\mathbf{S}(k^2)\) （\(k^2\) 的）（《代数学》第三章 §6 第 1 节 定义 1）。\(\mathfrak{sl}(2,k)\) 的元素唯一地开拓为 \(\mathbf{S}(k^2)\) 的导子，从而给 \(\mathbf{S}(k^2)\) 以 \(\mathfrak{sl}(2,k)\) -模结构（第一章 §3 第 2 节）。设 \(\mathrm{V}(m)\) 是 \(\mathbf{S}(k^2)\) 中次数为 \(m\) 的齐次元素之集。则 \(\mathrm{V}(m)\) 是 \(\mathfrak{sl}(2,k)\) -子模（\(\mathbf{S}(k^2)\) 的），维数为 \(m + 1\) ，即 \(m\) 次对称幂（\(\mathrm{V}(1) = k^2\) 的）（第三章附录）。若 \(m,n\) 是满足 \(0 \leq n \leq m\) 的整数，置

\[
e _ {n} ^ {(m)} = \binom {m} {n} u ^ {m - n} v ^ {n} \in \mathrm{V} (m).
\]

命题 3. 对任何整数 \(m \geq 0\) ，\(V(m)\) 是绝对单 \(\mathfrak{sl}(2, k)\) -模。在此模中，\(e_0^{(m)} = u^m\) 是权为 \(m\) 的本原元。

我们有 \(X_{+}u^{m}=0\) 与 \(Hu^{m}=mu^{m}\) ，故 \(u^{m}\) 是权 m 的本原元。\(V(m)\) 的由 \(u^{m}\) 生成的子模维数为 \(m+1\) （命题 2 (i)），故等于 \(V(m)\) 。由命题 2 (v)，\(V(m)\) 绝对单。

定理 1. 每个单 \(\mathfrak{sl}(2,k)\) -模（维数 \(n\) 有限）都同构于 \(\mathrm{V}(n - 1)\) 。每个有限维 \(\mathfrak{sl}(2,k)\) -模都是同构于诸模 \(\mathrm{V}(m)\) 的子模的直和。

这由引理 4 与命题 1、2、3 得出。

注. 1) \(\mathfrak{sl}(2,k)\) 的伴随表示在 \(\mathfrak{sl}(2,k)\) 上定义了单 \(\mathfrak{sl}(2,k)\) -模结构。此模同构于 V(2)，同构映射把 \(u^2\) 映到 \(X_+\) ，把 \(2uv\) 映到 \(-H\) ，把 \(v^2\) 映到 \(X_-\) 。

\begin{enumerate}
\def\labelenumi{\arabic{enumi})}
\setcounter{enumi}{1}
\tightlist
\item
  对 \(n \geq 0\) 且 m \textgreater{} n，
\end{enumerate}

\[
X _ {-} e _ {n} ^ {(m)} = - (m - n) \binom{m}{n} u ^ {m - n - 1} v ^ {n + 1} = - (n + 1) e _ {n + 1} ^ {(m)}.
\]

因此，\((e_0^{(m)}, e_1^{(m)}, \ldots, e_m^{(m)})\) 是 \(V(m)\) 的与本原元 \(e_0^{(m)}\) 相伴的基。

\begin{enumerate}
\def\labelenumi{\arabic{enumi})}
\setcounter{enumi}{2}
\tightlist
\item
  设 \(\Phi\) 是 \(V(m)\) 上的双线性型，使得
\end{enumerate}

\[
\begin{array}{c} \varPhi (e _ {n} ^ {(m)}, e _ {n ^ {\prime}} ^ {(m)}) = 0 \text {if} n + n ^ {\prime} \neq m \\ \varPhi (e _ {n} ^ {(m)}, e _ {m - n} ^ {(m)}) = (- 1) ^ {n} \binom {m} {n}. \end{array}
\]

若 \(x = au + bv\) 且 \(y = cu + dv\) ，则 \(\Phi(x^{m}, y^{m}) = (ad - bc)^{m}\) 。现在容易验证 \(\Phi\) 是不变的，且 \(\Phi\) 当 m 为偶时对称，当 m 为奇时交错。

命题 4. 设 E 是有限维 \(\mathfrak{sl}(2,k)\) -模，\(m\) 是整数 \(\geq 0\) ，\(\mathrm{P}_m\) 是权为 \(m\) 的本原元之集。设 L 是从 \(\mathfrak{sl}(2,k)\) -模 \(\mathrm{V}(m)\) 到 \(\mathfrak{sl}(2,k)\) -模 E 的同态组成的向量空间。映射 \(f \mapsto f(u^m)\) （从 L 到 E 的）是线性的、单射的，其像为 \(\mathrm{P}_m \cup \{0\}\) 。

此映射显然线性，且它是单射，因为 \(u^{m}\) 生成 \(\mathfrak{sl}(2,k)\) -模 \(\mathrm{V}(m)\) 。若 \(f \in L\) ，

\[
X _ {+} (f (u ^ {m})) = f (X _ {+} u ^ {m}) = 0, \quad H (f (u ^ {m})) = f (H u ^ {m}) = m f (u ^ {m})
\]

故 \(f(u^{m}) \in \mathrm{P}_{m} \cup \{0\}\) 。设 \(e \in \mathrm{P}_{m}\) ，\(V\) 是 \(E\) 的由 \(e\) 生成的子模。由命题 1，存在从模 \(V(m)\) 到模 \(V\) 的同构，把 \(u^{m}\) 映到 \(e\) 。于是 \(L(u^{m}) = P_{m} \cup \{0\}\) 。

推论. E 的型为 V(m) 的同构型分支的长度为

\(\dim (\mathrm{P}_m\cup \{0\} .\)

